\chapter*{Cómo leer este tutorial}
\addcontentsline{toc}{chapter}{Cómo leer este tutorial}
\markboth{Cómo leer este tutorial}{}

Este documento recorre Alexandria de punta a punta con un único ejemplo que se parece a un problema de
verdad: una distribuidora eléctrica que comparte curvas de carga cuartohorarias con comercializadoras,
agregadores de flexibilidad y un grupo de investigación. No hay colores ni acciones de juguete: hay finalidades
de uso, agregación mínima, plazos de retención y jurisdicciones, que es de lo que tratan las políticas reales.

El recorrido tiene cinco partes y cada una se apoya en la anterior:

\begin{enumerate}[leftmargin=*, itemsep=2pt]
  \item \textbf{El modelo mental.} Qué problema resuelve Alexandria, qué hace cada binario y cuál es el caso de estudio.
  \item \textbf{El lenguaje de perfiles.} Cómo se escribe un \cod{OdrlProfile} y un \cod{SpaceProfile}, qué genera
        \cod{alexandriac}, cómo valida y cómo se versiona.
  \item \textbf{El hub.} Cómo \cod{alexandria-server} convierte ese YAML en una versión firmada, con fechas y auditoría.
  \item \textbf{El borde.} Cómo \cod{alexandria-client} sincroniza esas versiones y decide sobre cada mensaje DSP.
  \item \textbf{Revisión de SHACL.} Qué está hecho, qué no, qué falla cuando se ejecutan de verdad las \emph{shapes}
        generadas y qué haría falta para cerrarlo.
\end{enumerate}

\section*{Convenciones}

\noindent\begin{tabularx}{\linewidth}{@{}L{3.6cm}Y@{}}
\toprule
\textbf{Aspecto} & \textbf{Significado} \\
\midrule
Fondo azul claro & Comando que se escribe en la terminal. \\
Fondo blanco, letra pequeña & Salida real del comando, copiada tal cual (recortada con \texttt{…} cuando es larga). \\
Fondo gris & Contenido de un fichero: YAML, JSON o Turtle. \\
Caja azul, «Idea clave» & El concepto que conviene retener de la sección. \\
Caja ámbar, «Trampa habitual» & Un comportamiento que sorprende la primera vez. \\
Caja gris, «Bajo el capó» & Dónde está en el código lo que se acaba de ver. \\
Caja roja, «Hallazgo» & Un defecto o una carencia detectada durante la revisión. \\
\bottomrule
\end{tabularx}

\section*{Qué necesitas}

\begin{itemize}[leftmargin=*, itemsep=2pt]
  \item Go 1.22 o posterior y \cod{make} para compilar (\texttt{make build}).
  \item \cod{curl} y \cod{python3} (solo biblioteca estándar) para las llamadas y para leer JSON.
  \item Para reproducir la revisión de SHACL del capítulo~\ref{ch:shacl}: un entorno virtual de Python con
        \cod{pyshacl} y \cod{jsonschema}.
\end{itemize}

\section*{Los ficheros del tutorial}

Todo lo que se cita vive en \cod{tutorial/manual/examples/}:

\noindent\begin{tabularx}{\linewidth}{@{}L{5.4cm}Y@{}}
\toprule
\textbf{Ruta} & \textbf{Contenido} \\
\midrule
\cod{dsl/metering-1.0.0.yaml} & El perfil ODRL de medida, solo. \\
\cod{dsl/flexgrid-1.0.0.yaml} & Perfil 1.0.0 más el \cod{SpaceProfile} 1.0.0: el primer borrador que se publica. \\
\cod{dsl/flexgrid-1.1.0.yaml} & Versión compatible: una finalidad nueva y una acción obsoleta. \\
\cod{dsl/flexgrid-2.0.0.yaml} & Versión de ruptura: la acción obsoleta se retira. \\
\cod{dsl/health-0.3.0.yaml} & Segundo perfil, de uso secundario de datos de salud. \\
\cod{vendor/dsp/2025-1/} & Copia local reducida de DSP 2025-1 para el \emph{binding} vendorizado. \\
\cod{msgs/} & Los mensajes DSP que se envían: \cod{m1}–\cod{m6}, el acuerdo \cod{a1} y las transferencias. \\
\cod{rules/} & Las tres reglas del cliente que se adaptan. \\
\cod{lab/} & \cod{server.yaml}, \cod{client.yaml} y \cod{env.sh}, el guion del laboratorio. \\
\cod{shacl/harness.py} & El banco de pruebas de la revisión de SHACL. \\
\bottomrule
\end{tabularx}

\begin{notaPractica}[Sobre los puertos]
El laboratorio de este tutorial escucha en \cod{18443} (hub), \cod{19400} (IdP), \cod{18282} (conector),
\cod{19080}/\cod{19081} (tráfico) y \cod{19090} (administración). Son distintos de los de \cod{config/server.yaml}
y \cod{config/client.yaml} a propósito: así puede convivir con el laboratorio por defecto, que usa otra base de
datos, otro compilador de borradores y otro \emph{namespace}.
\end{notaPractica}
